\section{Selected exact finite determinations}
\label{sec:selected-exact}

The all-\(n\) results above are accompanied by exact finite determinations whose
proofs combine a structural reduction with a complete certificate.  This
section records the strongest examples in enough detail to expose the
exhaustiveness argument.

\subsection{Cutting centers of trees: \texorpdfstring{\seq{A002887}}{A002887}}

For a tree \(T\) of order \(N\), delete a vertex \(v\) and let the component
orders be \(s_1,\ldots,s_d\).  The cutting number is
\begin{equation}
 c_T(v)=\sum_{i<j}s_is_j
       =\frac{(N-1)^2-\sum_i s_i^2}{2}.
 \label{eq:a002887-cut}
\end{equation}
Harary and Ostrand proved that the vertices of maximum cutting number form a
path \citep{HararyOstrand1971}.  Suppose the intended center path is
\(v_1,\ldots,v_k\).  Removing its path edges leaves rooted components of orders
\(b_1,\ldots,b_k\).  Let \(q_i\) be the sum of the squared orders of the
off-path branches incident with \(v_i\), and put
\[
 L_i=\sum_{j<i}b_j,
 \qquad
 R_i=\sum_{j>i}b_j.
\]
The path vertices have a common cutting number exactly when
\begin{equation}
 L_i^2+R_i^2+q_i=K
 \qquad(1\le i\le k)
 \label{eq:a002887-profile}
\end{equation}
for one integer \(K\).  Their common value is
\(( (N-1)^2-K)/2\).

For fixed \(N\) and proposed center value \(C\), call a rooted tree of order
\(s\) safe when every vertex in it has cutting number strictly below \(C\)
after the root is attached to an outside component of order \(N-s\).  If the
root has safe child-subtrees of orders \(t_1,\ldots,t_d\), with
\(\sum t_i=s-1\) and \(q=\sum t_i^2\), its cutting number is
\begin{equation}
 (N-s)(s-1)+\frac{(s-1)^2-q}{2}.
 \label{eq:a002887-safe}
\end{equation}
Thus safety is characterized recursively by an exact integer-partition
dynamic program.  This condition is necessary by deleting the root and
sufficient by attaching witnessing safe subtrees.  Enumerating positive
compositions \((b_1,\ldots,b_k)\), exact attainable square sums in
\cref{eq:a002887-profile}, and the safe-subtree spectra therefore loses no
trees.

Searching orders increasingly gives
\begin{center}
\begin{tabular}{@{}c*{10}{r}@{}}
\toprule
\(k\) & 1 & 2 & 3 & 4 & 5 & 6 & 7 & 8 & 9 & 10 \\
\midrule
minimum \(N\) & 3 & 4 & 7 & 10 & 48 & 48 & 122 & 122 & 264 & 264 \\
\bottomrule
\end{tabular}
\end{center}
The value at \(k=5\) corrects 50 to 48.  Every accepted profile is realized as
an explicit edge list; an independent checker deletes every vertex,
recomputes all component orders and cutting numbers, and verifies that the
maximizers are exactly the intended path.

\subsection{Three square and three cube representations: \texorpdfstring{\seq{A273354}}{A273354}}

Let \(r_2^+(N)\) and \(r_3^+(N)\) be the numbers of unordered positive pairs
representing \(N\) as a sum of two squares and two cubes.  The exact third term
is
\begin{equation}
 \boxed{\seq{A273354}(3)=11177126654841000000.}
 \label{eq:a273354-value}
\end{equation}
Put \(N_0=11177126654841000000=3343221000^2\).  Direct integer arithmetic gives
exactly the three cube representations
\begin{align*}
N_0&=279300^3+2234400^3\\
   &=790020^3+2202480^3\\
   &=1256850^3+2094750^3.
\end{align*}
Its factorization is
\[
 N_0=2^6 3^6 5^6 7^6 19^4.
\]
The sum-of-two-squares divisor formula gives exactly three positive unordered
square representations; the certificate lists them explicitly.

For minimality, suppose \(M<N_0\) has at least three positive-cube
representations.  Dividing all six coordinates by their common gcd writes
\(M=g^3m\), where \(m\) is a primitive three-cube value.  Thus every candidate
occurs among the cubic multiples below \(N_0\) of the primitive values in the
complete \seq{A003825} ancillary list.  There are 110969 distinct candidates.
Each was factored exactly using deterministic Miller-Rabin and Pollard-rho
arithmetic, and the positive-square-pair formula was evaluated from the prime
factorization.  None has \(r_2^+(M)=3\).  An independent verifier regenerates
the candidate set, checks primality and factor products, recomputes the square
counts, and directly scans all cube pairs of \(N_0\).

\subsection{Missing polygonal cases: \texorpdfstring{\seq{A363253}}{A363253}}

For \(s\ge3\), let
\[
 P_s(k)=\frac{(s-2)k^2-(s-4)k}{2}
\]
be the \(k\)-th positive \(s\)-gonal number.  Representations are sums of
distinct positive \(s\)-gonal numbers, including the singleton representation.
The missing cases are
\begin{equation}
 \boxed{\seq{A363253}(6)=\seq{A363253}(7)=-1.}
 \label{eq:a363253-values}
\end{equation}

The proof uses a conductor lemma.  If subset sums of
\(P_s(1),\ldots,P_s(m)\) cover every integer in \([L,U]\) and
\(P_s(m+1)\le U-L+1\), then adjoining \(P_s(m+1)\) makes the old and shifted
intervals overlap.  Since the successive polygonal numbers in the relevant
range grow by a factor less than two, the overlap persists indefinitely.
Exact bitset calculations show that
\begin{align*}
 P_6(1),\ldots,P_6(11)&\text{ cover }[268,678],\\
 P_7(1),\ldots,P_7(13)&\text{ cover }[388,1523].
\end{align*}
Therefore every integer at least 268 is a distinct hexagonal subset sum and
every integer at least 388 is a distinct heptagonal subset sum.

For \(k\ge68\), the six differences
\(P_6(k)-P_6(k-d)\), \(1\le d\le6\), lie beyond the conductor and below
\(P_6(k-d)\).  Appending the distinct largest summands
\(P_6(k-d)\) gives six non-singleton representations, hence at least seven in
total.  The analogous seven differences for heptagonal numbers work for
\(k\ge79\), giving at least eight representations.  Exact finite-prefix subset
DP rules out the desired multiplicities before these thresholds.

\subsection{Minimally 2-edge-connected graphs: \texorpdfstring{\seq{A001072}}{A001072}}

The exact new term is
\begin{equation}
 \boxed{\seq{A001072}(15)=41142.}
 \label{eq:a001072-value}
\end{equation}
By the Whitney ear theorem \citep{Whitney1932}, every 2-edge-connected graph is
built from a cycle by adding ears.  In a minimal 2-edge-connected simple graph,
every later ear has a new internal vertex, and every prefix remains minimal.
The generator recursively adds all open ears with at least one new internal
vertex and all simple closed ears with at least two, retaining exactly the
graphs that are connected and bridgeless but cease to be so after any edge
deletion.  Canonical graph6 records produced by nauty
\citep{McKayPiperno2014} remove isomorphic duplicates.

The computation gives 203421 accepted precanonical candidates and 41142
canonical records at order 15, after reproducing every published term through
order 14.  An independent verifier uses a different minimality test: for each
edge \(e\), it checks that \(G-e\) is connected and finds an edge \(f\) for
which \(G-\{e,f\}\) is disconnected.  A separate canonicalization pass leaves
the count and record hash unchanged.

\subsection{Free-monoid multiplication programs: \texorpdfstring{\seq{A075099}}{A075099}}

Every binary word of length \(n\) requires one final multiplication.  A
shorter word contributes one additional multiplication exactly when it is
retained as a reusable intermediate.  Binary variables select these
intermediates.  For every selected word and every required target, split
variables require a cut whose two factors are generators or selected shorter
words.  Since factors are strictly shorter, the local split conditions are
equivalent to a valid multiplication program ordered by word length.

The resulting mixed-integer program minimizes \(2^n\) plus the number of
selected intermediates.  Completed HiGHS runs have matching primal and global
bounds and zero gap, giving
\begin{equation}
 \boxed{a(7)=151,\quad a(8)=276,\quad a(9)=556,\quad
        a(10)=1066,\quad a(12)=4171.}
 \label{eq:a075099-values}
\end{equation}
The selected programs are checked independently by recursively verifying a
valid factorization of every intermediate and every target.  The stored
time-limited run at \(n=11\) proves only
\(2118\le a(11)\le2138\); it is not used in
\cref{eq:a075099-values}.

\subsection{Triangular-lattice labeling: \texorpdfstring{\seq{A337433}}{A337433}}

For the 28 nodes of the order-seven triangular array, binary variables assign
one label in \(\{1,\ldots,7\}\) to each node.  If a node receives label \(k\),
the model requires an occurrence of every label \(1,\ldots,k-1\) among its
neighbors.  This is exactly the defining local condition.  A completed
53,115-node branch-and-bound computation has objective and global bound 87
with zero gap.  The resulting labeling is independently checked node by node,
and therefore
\begin{equation}
 \boxed{\seq{A337433}(7)=87.}
 \label{eq:a337433-value}
\end{equation}

\subsection{Additional exact values}

\begin{table}[htbp]
\centering
\small
\begin{tabularx}{\textwidth}{@{}l l X@{}}
\toprule
Entry & Exact result & Certificate \\
\midrule
A323560 & \(a(19)=36001752\) & Exhaustive reversed trapped-knight paths, with a Held-Karp lower bound for the unvisited required neighbors. \\
A368353 & \(a(6)=-2579770\) & \multirow{3}{=}{Exhaustive \(11!/2\) reversal classes of \(6\times6\) Hankel matrices; exact fraction-free Bareiss determinants \citep{Bareiss1968}.} \\
A368354 & \(a(6)=2256911\) & \\
A368355 & \(a(6)=2579770\) & \\
\bottomrule
\end{tabularx}
\caption{Further exact finite determinations.}
\label{tab:additional-exact}
\end{table}
